% GENERATED FILE. Edit canonical lessons, then run npm run generate:books.
% canonical-source-hash: fnv1a64:24d0b807c0d3b1f2
% canonical-lessons: ML-C269-asvasippikkuka, ML-C269-prartthikkuka, ML-C269-alankarikkuka, ML-C269-janikkuka, ML-C269-marikkuka

\chapter{To comfort, To pray, To decorate, To be born, To die}
\label{ch:ml-to-comfort-to-pray-to-decorate-to-be-born-to-die}
\hlchaptermodality{269}

\begin{chapteropening}
\textbf{By the end of this chapter:} \emph{I can say to comfort, to pray, to decorate, to be born and to die.}

Complete the last lesson of chapter 269: I can say to comfort, to pray, to decorate, to be born and to die.
\end{chapteropening}

\section[\texorpdfstring{\ml{ആശ്വസിപ്പിക്കുക} (āśvasippikkuka)}{ആശ്വസിപ്പിക്കുക (āśvasippikkuka)}]{\ml{ആശ്വസിപ്പിക്കുക} (āśvasippikkuka) — to comfort, to console}

\label{lesson:ML-C269-asvasippikkuka}



\begin{quote}
Before the new one: say the Malayalam for to register, then the Malayalam for to deliver.
\end{quote}

\subsection*{You'll want to know: \ml{ആശ്വസിപ്പിക്കുക}}
\textbf{\ml{ആശ്വസിപ്പിക്കുക}} — \emph{āśvasippikkuka} — \textquotedblleft{}to comfort, to console\textquotedblright{}.

\begin{grammarlens}[title={What you've built}]
One more word for this chapter.
\end{grammarlens}

\subsection*{Your turn}
\begin{itemize}
\raggedright
  \item \emph{Say it:} \emph{āśvasippikkuka}
  \item \emph{Say it:} \emph{āśvasippikkuka}, once more
  \item \emph{Say it:} say \emph{āśvasippikkuka}
  \item \emph{Recall:} say the Malayalam for to register, then the Malayalam for to deliver, then say \emph{āśvasippikkuka} again
\end{itemize}

\begin{tcolorbox}[breakable,colback=teal!4,colframe=teal!35!black,title={Before you move on}]
What does \ml{ആശ്വസിപ്പിക്കുക} mean? (\textquotedblleft{}To comfort, to console\textquotedblright{}.) Say it once more.
\end{tcolorbox}
\section[\texorpdfstring{\ml{പ്രാർത്ഥിക്കുക} (prārtthikkuka)}{പ്രാർത്ഥിക്കുക (prārtthikkuka)}]{\ml{പ്രാർത്ഥിക്കുക} (prārtthikkuka) — to pray}

\label{lesson:ML-C269-prartthikkuka}



\begin{quote}
Before the new one: say the Malayalam for to deliver, then the Malayalam for to comfort.
\end{quote}

\subsection*{You'll want to know: \ml{പ്രാർത്ഥിക്കുക}}
\textbf{\ml{പ്രാർത്ഥിക്കുക}} — \emph{prārtthikkuka} — \textquotedblleft{}to pray\textquotedblright{}.

\begin{grammarlens}[title={What you've built}]
One more word for this chapter.
\end{grammarlens}

\subsection*{Your turn}
\begin{itemize}
\raggedright
  \item \emph{Say it:} \emph{prārtthikkuka}
  \item \emph{Say it:} \emph{prārtthikkuka}, once more
  \item \emph{Say it:} say \emph{prārtthikkuka}
  \item \emph{Recall:} say the Malayalam for to deliver, then the Malayalam for to comfort, then say \emph{prārtthikkuka} again
\end{itemize}

\begin{tcolorbox}[breakable,colback=teal!4,colframe=teal!35!black,title={Before you move on}]
What does \ml{പ്രാർത്ഥിക്കുക} mean? (\textquotedblleft{}To pray\textquotedblright{}.) Say it once more.
\end{tcolorbox}
\section[\texorpdfstring{\ml{അലങ്കരിക്കുക} (alaṅkarikkuka)}{അലങ്കരിക്കുക (alaṅkarikkuka)}]{\ml{അലങ്കരിക്കുക} (alaṅkarikkuka) — to decorate}

\label{lesson:ML-C269-alankarikkuka}



\begin{quote}
Before the new one: say the Malayalam for to comfort, then the Malayalam for to pray.
\end{quote}

\subsection*{You'll want to know: \ml{അലങ്കരിക്കുക}}
\textbf{\ml{അലങ്കരിക്കുക}} — \emph{alaṅkarikkuka} — \textquotedblleft{}to decorate\textquotedblright{}.

\begin{grammarlens}[title={What you've built}]
One more word for this chapter.
\end{grammarlens}

\subsection*{Your turn}
\begin{itemize}
\raggedright
  \item \emph{Say it:} \emph{alaṅkarikkuka}
  \item \emph{Say it:} \emph{alaṅkarikkuka}, once more
  \item \emph{Say it:} say \emph{alaṅkarikkuka}
  \item \emph{Recall:} say the Malayalam for to comfort, then the Malayalam for to pray, then say \emph{alaṅkarikkuka} again
\end{itemize}

\begin{tcolorbox}[breakable,colback=teal!4,colframe=teal!35!black,title={Before you move on}]
What does \ml{അലങ്കരിക്കുക} mean? (\textquotedblleft{}To decorate\textquotedblright{}.) Say it once more.
\end{tcolorbox}
\section[\texorpdfstring{\ml{ജനിക്കുക} (janikkuka)}{ജനിക്കുക (janikkuka)}]{\ml{ജനിക്കുക} (janikkuka) — to be born}

\label{lesson:ML-C269-janikkuka}



\begin{quote}
Before the new one: say the Malayalam for to pray, then the Malayalam for to decorate.
\end{quote}

\subsection*{You'll want to know: \ml{ജനിക്കുക}}
\textbf{\ml{ജനിക്കുക}} — \emph{janikkuka} — \textquotedblleft{}to be born\textquotedblright{}.

\begin{grammarlens}[title={What you've built}]
One more word for this chapter.
\end{grammarlens}

\subsection*{Your turn}
\begin{itemize}
\raggedright
  \item \emph{Say it:} \emph{janikkuka}
  \item \emph{Say it:} \emph{janikkuka}, once more
  \item \emph{Say it:} say \emph{janikkuka}
  \item \emph{Recall:} say the Malayalam for to pray, then the Malayalam for to decorate, then say \emph{janikkuka} again
\end{itemize}

\begin{tcolorbox}[breakable,colback=teal!4,colframe=teal!35!black,title={Before you move on}]
What does \ml{ജനിക്കുക} mean? (\textquotedblleft{}To be born\textquotedblright{}.) Say it once more.
\end{tcolorbox}
\section[\texorpdfstring{\ml{മരിക്കുക} (marikkuka)}{മരിക്കുക (marikkuka)}]{\ml{മരിക്കുക} (marikkuka) — to die}

\label{lesson:ML-C269-marikkuka}



\begin{quote}
Before the new one: say the Malayalam for to decorate, then the Malayalam for to be born.
\end{quote}

\subsection*{You'll want to know: \ml{മരിക്കുക}}
\textbf{\ml{മരിക്കുക}} — \emph{marikkuka} — \textquotedblleft{}to die\textquotedblright{}.

\begin{grammarlens}[title={What you've built}]
One more word for this chapter.
\end{grammarlens}

\subsection*{Your turn}
\begin{itemize}
\raggedright
  \item \emph{Say it:} \emph{marikkuka}
  \item \emph{Say it:} \emph{marikkuka}, once more
  \item \emph{Say it:} say \emph{marikkuka}
  \item \emph{Recall:} say the Malayalam for to decorate, then the Malayalam for to be born, then say \emph{marikkuka} again
\end{itemize}

\begin{tcolorbox}[breakable,colback=teal!4,colframe=teal!35!black,title={Before you move on}]
What does \ml{മരിക്കുക} mean? (\textquotedblleft{}To die\textquotedblright{}.) Say it once more.
\end{tcolorbox}
